%!TEX root = ../template.tex
\chapter{Notas de revisão das listas preliminares}
\label{app:revisao-editorial}

Estas notas dizem respeito aos ficheiros de siglas, símbolos e glossário.
Este apêndice só é incluído quando os comentários de revisão estão ativos.

% [REV-2026-09-13][FR-01][P3][APRESENTACAO]
\revisaotese{FR-01}{P3}{APRESENTACAO}
{Substituir as siglas de exemplo pelas usadas na tese}
{Este ficheiro ainda contém abbrev, aaa, xpto e outras entradas do
template. Criar apenas as necessárias: CDU, SGE, USE, EnPI, EnB, GNE,
EII, OLS, HAC, MBB, MAE, CUSUM, PRISMA-ScR, entre outras realmente usadas.
Definir equivalências EnPI/IDE e EnB/CER uma vez e escolher a notação
principal. Uniformizar kt/d, t GNE/d e grafia fuel gás/fuel gas.}%
% [/REV-2026-09-13]
% [REV-2026-09-13][FR-02][P3][NOTACAO]
\revisaotese{FR-02}{P3}{NOTACAO}
{Definir a notação física e estatística com unidades}
{Trocar símbolos de exemplo por E\_t, Q\_t, Z\_t, a, b, e\_t, sigma residual,
sigma da inovação, rho, delta\_MAE, cobertura, k e h do CUSUM.
Dar unidades e separar índice temporal de ano. A unidade do declive
com E em t GNE/d e Q em kt/d é t GNE/kt: o dia cancela.
Distinguir IC de parâmetro/média, intervalo preditivo e banda de alarme
no apêndice de matemática e nas legendas.}%
% [/REV-2026-09-13]
% [REV-2026-09-13][FR-03][P3][APRESENTACAO]
\revisaotese{FR-03}{P3}{APRESENTACAO}
{Remover conteúdo demonstrativo e fixar termos operacionais}
{O glossário contém ainda a entrada computer e uma citação de exemplo
Artho04 sem relação com a tese. Substituir por definições curtas de
baseline, normalização, referência congelada, fator estático, desvio,
excursão, episódio de alarme, calibração e reconciliação de fontes,
se o glossário for mantido. As distinções resolvem ambiguidades
recorrentes no texto, não são apenas limpeza do template.}%
% [/REV-2026-09-13]
